\section*{Chapitre \ref{ch:b3:dna-repair} --- Stabilité du génome~: lésions de l'ADN, réparation et recombinaison}

\begin{solution}{exo:b3:dna-repair:1}
Uracile~: spontanée (\omterm{def:b3:dna-repair:lesions}{désamination} d'un C), \omterm{def:b3:dna-repair:excision}{réparation par excision de
base} grâce à l'uracile-ADN-glycosylase. Dimère de thymines~:
environnementale (UV), \omterm{def:b3:dna-repair:excision}{réparation par excision de nucléotides}.
\omterm{def:b3:dna-repair:lesions}{8-oxoguanine}~: spontanée (oxydation), \omterm{def:b3:dna-repair:excision}{réparation par excision de base}
grâce à la \omterm{def:b3:dna-repair:excision}{glycosylase} de la 8-oxoG. Mésappariement G--T derrière la
fourche~: erreur de réplication, \omterm{prop:b3:dna-repair:fidelity}{réparation des mésappariements}.
\omterm{def:b3:dna-repair:lesions}{Cassure double brin} en G1~: environnementale (rayonnement) ou
spontanée~; \omterm{def:b3:dna-repair:dsb}{jonction d'extrémités non homologues}, puisqu'aucune
chromatide sœur n'est disponible.
\end{solution}

\begin{solution}{exo:b3:dna-repair:2}
L'excision de nucléotides reconnaît une propriété physique commune à
toutes les \omterm{def:b3:dna-repair:lesions}{lésions} volumineuses --- la déformation de l'hélice, ou une
ARN polymérase bloquée --- et retire une rustine contenant ce qui l'a
causée. Les petites \omterm{def:b3:dna-repair:lesions}{lésions} (uracile, 8-oxoG, bases alkylées) ne
déforment guère l'hélice, il faut donc les reconnaître par leur
chimie~: une \omterm{def:b3:dna-repair:excision}{glycosylase} par type de base fautive.
\end{solution}

\begin{solution}{exo:b3:dna-repair:3}
Un mésappariement est fait de deux bases normales, et seule celle du
brin neuf est fautive~; le système doit savoir quel brin est neuf.
\emph{E.~coli} marque le brin parental de groupements méthyle aux
sites GATC, que le brin neuf n'a pas pendant quelques minutes~; MutH
entaille le brin non méthylé. Chez un mutant \emph{dam} aucun brin
n'est méthylé~: MutH entaille l'un ou l'autre, la réparation retire la
bonne base aussi souvent que la mauvaise --- fixant la moitié des
erreurs en mutations (un mutateur) --- et des entailles sur les deux
brins à des sites GATC voisins créent des \omterm{def:b3:dna-repair:lesions}{cassures double brin}.
\end{solution}

\begin{solution}{exo:b3:dna-repair:4}
Environ $35\times 2 = 70$ \omterm{def:b3:dna-repair:lesions}{cassures double brin}. Au bout d'une heure la
plupart sont jointes~: avec une demi-vie voisine de \qty{30}{min},
environ $70/4 \approx 18$ foyers. Au bout de six heures presque toutes
ont été réparées ($70/2^{12} < 1$), ne laissant que quelques cassures
persistantes et complexes.
\end{solution}

\begin{solution}{exo:b3:dna-repair:5}
(a) $10^{-5}\times 10^{-2}\times 10^{-3} = 10^{-10}$~; $6.4\times
10^{9}\times 10^{-10} = 0.64$ mutation par division. (b) Sans
\omterm{prop:b3:dna-repair:fidelity}{réparation des mésappariements} $10^{-7}$~; $640$ par division. (c) Sans
relecture $10^{-8}$~; $64$ par division. Le mutant de la \omterm{prop:b3:dna-repair:fidelity}{réparation des
mésappariements} est le plus dangereux, mille fois au-dessus de la
normale.
\end{solution}

\begin{solution}{exo:b3:dna-repair:6}
$k = \ln 2/\qty{0.25}{h} = \qty{2.8}{h^{-1}}$~; $D = 10^{4}/24 =
\qty{420}{h^{-1}}$~; $L^{*} = 420/2.8 \approx 150$ sites abasiques
présents à un instant donné. Réparation dix fois plus lente~:
$L^{*} \approx \num{1500}$. La fourche passe une fois par chaque locus et y
rencontre les \omterm{def:b3:dna-repair:lesions}{lésions} présentes à cet instant~; moyenné sur le génome
elle en rencontre environ $L^{*}$ --- quelque $150$ normalement,
$\num{1500}$ sous le médicament.
\end{solution}

\begin{solution}{exo:b3:dna-repair:7}
$k_{\text{NHEJ}} = \ln 2/0.5 = \qty{1.39}{h^{-1}}$, $k_{\text{HR}} =
\ln 2/4 = \qty{0.17}{h^{-1}}$~; fraction HR $0.17/1.56 = 0.11$~;
demi-vie $\ln 2/1.56 = \qty{27}{min}$. Une fourche effondrée donne une
cassure à \emph{une seule extrémité}~: il n'y a pas de seconde
extrémité à joindre pour la NHEJ, seule la recombinaison avec la sœur
peut rétablir la fourche --- et c'est de toute façon la seule voie
fidèle.
\end{solution}

\begin{solution}{exo:b3:dna-repair:8}
(a) Lynch~: microsatellites instables (les longueurs diffèrent d'une
cellule tumorale à l'autre), cancers colorectal, de l'endomètre,
gastrique et de l'ovaire, réponse normale au soleil. (b) XP-A~:
microsatellites stables, coup de soleil extrême à la moindre
exposition, cancers de la peau et de l'œil dès l'enfance, et dans ce
groupe une dégénérescence neurologique progressive. (c) XP-V~:
microsatellites stables, excision normale donc coup de soleil plus
modéré, mais les dimères qui échappent à l'excision sont contournés par
des polymérases infidèles --- cancers cutanés, plus tardifs que dans le
groupe~A.
\end{solution}

\begin{solution}{exo:b3:dna-repair:9}
Cellules hépatiques~: Cre a excisé l'exon~3 des deux allèles (sous
forme d'un cercle qui se perd), donc le gène est inactif dans chaque
hépatocyte dès la naissance --- une invalidation propre au foie.
Neurones~: pas de Cre, les allèles floxés sont intacts et
fonctionnels. Orientations opposées~: Cre inverse l'exon et l'inverse
de nouveau indéfiniment, de sorte qu'à tout instant la moitié environ
des cellules porte l'exon inversé, non fonctionnel --- une mosaïque, et
non une délétion nette.
\end{solution}

\begin{solution}{exo:b3:dna-repair:10}
À mesure que LexA est clivé, sa concentration baisse progressivement.
Les opérateurs qui fixent LexA faiblement se libèrent dès une faible
baisse --- \emph{uvrA}, \emph{uvrB}, \emph{recA} --- donc la réparation
fidèle commence la première~; les opérateurs de \emph{sulA} et des
polymérases infidèles fixent LexA fortement et ne se libèrent que
lorsqu'il a presque disparu, c'est-à-dire lorsque les dégâts persistent
depuis longtemps. Une bactérie sans polymérases infidèles mourrait aux
fourches bloquées par des \omterm{def:b3:dna-repair:lesions}{lésions} qu'elle ne saurait contourner et,
faute de mutagenèse induite, acquerrait bien plus lentement la
résistance à l'antibiotique --- désavantage pour la bactérie, avantage
pour le patient~; des inhibiteurs de la \omterm{prop:b3:dna-repair:sos}{réponse SOS} sont étudiés pour
cette raison.
\end{solution}

\begin{solution}{exo:b3:dna-repair:11}
Chaîne~: PARP inhibée $\to$ les cassures simple brin persistent $\to$
une fourche de réplication transforme chacune en une \omterm{def:b3:dna-repair:lesions}{cassure double
brin} à une seule extrémité $\to$ la réparation exige la \omterm{def:b3:dna-repair:dsb}{recombinaison
homologue} $\to$ la cellule tumorale, privée de \omterm{def:b3:dna-repair:dsb}{BRCA}, ne le peut pas
$\to$ mauvaises jonctions, pertes de chromosomes, mort. Les cellules
normales, avec un allèle fonctionnel, recombinent et survivent.
Résistance~: (1) une seconde mutation de \emph{\omterm{def:b3:dna-repair:dsb}{BRCA}} qui rétablit le
cadre de lecture et une protéine fonctionnelle --- se détecte par
séquençage de la tumeur ou de l'ADN tumoral circulant dans le sang, et
par le retour des foyers de \omterm{def:b3:dna-repair:dsb}{Rad51} après irradiation~; (2) la perte des
protéines qui bloquent la résection (53BP1--shieldine), qui rétablit en
partie la recombinaison sans BRCA1 --- se détecte par leur absence au
marquage et par le retour des foyers de \omterm{def:b3:dna-repair:dsb}{Rad51}~; également des mutations
de PARP1 qui empêchent son piégeage sur l'ADN, ou des pompes d'efflux,
détectées respectivement par séquençage et par expression.
\end{solution}

\begin{solution}{exo:b3:dna-repair:12}
Deux duplex homologues, l'un porteur de $A$ et l'autre de $a$,
échangent des brins isolés en travers du marqueur~: chacun a désormais
un brin de chaque parent sur ce tronçon --- un hétéroduplex portant un
mésappariement $A/a$. Si la \omterm{prop:b3:dna-repair:fidelity}{réparation des mésappariements} convertit le
mésappariement du duplex envahi en $A$, trois chromatides portent $A$
et une $a$~: $3{:}1$~; convertie en $a$~: $1{:}3$~; laissée non
réparée, la chromatide ségrège $A$ et $a$ à la première mitose après la
méiose et donne deux spores filles différentes (ségrégation
post-méiotique, $5{:}3$ dans un asque à huit spores). Le tronçon
d'hétéroduplex ne fait que quelques centaines à quelques milliers de
paires de bases autour de la cassure initiatrice, de sorte que seuls
les marqueurs qui s'y trouvent peuvent être convertis --- et ils sont,
par construction, voisins du crossing-over.
\end{solution}

\begin{solution}{pb:b3:dna-repair:1}
\textbf{1.} $10^{4} + 400 + 10^{4} \approx 2\times 10^{4}$ \omterm{def:b3:dna-repair:lesions}{lésions} par
jour, $7\times 10^{6}$ par an --- à peu près une par kilobase de génome
et par an.
\textbf{2.} $10^{-5}\times 10^{-2}\times 10^{-3} = 10^{-10}$ par bp~;
$6.4\times 10^{9}\times 10^{-10} = 0.64$ mutation par division.
\textbf{3.} $0.64\times 0.015 \approx 0.01$ mutation codante par
division~; sur $50$ divisions, environ $0.5$.
\textbf{4.} $10^{-7}$ par bp, $640$ mutations par division. Cellules
glissées~: sans réparation
$10^{-4}\times 40\times 10^{9} = 4\times 10^{6}$~; avec elle
$10^{-7}\times 40\times 10^{9} = 4000$ --- un facteur mille, visible
sous forme d'\omterm{def:b3:dna-repair:mmr}{instabilité des microsatellites}.
\textbf{5.} L'uracile est étranger à l'ADN, il est reconnu par
l'uracile-ADN-glycosylase et excisé~: réparé. La thymine issue d'une
5-méthylcytosine est une base normale en face d'un G~; rien ne la
signale comme fautive hors réplication, et elle devient une mutation
C$\to$T --- le mécanisme qui a appauvri le génome en \omterm{def:b3:chromatin-epigenetics:methylation}{CpG}.
\textbf{6.} La \omterm{def:b3:dna-repair:excision}{glycosylase} de la 8-oxoG retire la base oxydée tant
qu'elle est encore en face d'un C (avant la réplication)~; la seconde
\omterm{def:b3:dna-repair:excision}{glycosylase} retire un A mal incorporé en face d'une 8-oxoG (après la
réplication, ce qui rend une chance de réparer le G)~; l'hydrolase
détruit le dGTP oxydé pour qu'il ne soit pas incorporé en face d'un A.
Lorsque les trois échouent, la 8-oxoG s'apparie à A et le tour suivant
fixe une transversion G$\cdot$C$\to$T$\cdot$A.
\textbf{7.} Normale $k = \ln 2/2 = \qty{0.35}{h^{-1}}$~; XP $k = \ln 2/70
= \qty{0.0099}{h^{-1}}$.
\textbf{8.} Normale $10^{5} e^{-2.77} \approx 6200$~; XP $10^{5}
e^{-0.079} \approx \num{92000}$.
\textbf{9.} Mutations $6.2$ et $92$ par cellule~; codantes $0.09$ et $1.4$.
\textbf{10.} $500$ expositions~: $47$ mutations codantes (normale) et
$690$ (XP), contre $0.5$ dues à la réplication --- le soleil domine la
charge mutationnelle de la peau même chez une personne normale.
\textbf{11.} Séquence codante $9.6\times 10^{7}$ bp~; les gènes moteurs
en représentent $4\times 10^{4}/9.6\times 10^{7} = 4.2\times 10^{-4}$.
Moyenne des atteintes $690\times 4.2\times 10^{-4} = 0.29$~;
$P(\ge 1) = 1 - e^{-0.29} =
0.25$. Sur $10^{9}$ cellules, $2.5\times 10^{8}$ portent une mutation
motrice (enfant normal~: moyenne $0.02$, \qty{2}{\%}).
\textbf{12.} Un dimère n'est pas une mutation~; il ne le devient que
lorsqu'une fourche le recopie par contournement infidèle, de sorte que
ce sont les cellules qui se divisent avec beaucoup de dimères non
réparés qui mutent. L'ultraviolet n'atteint que la peau et l'œil~; les
tissus internes ne reçoivent aucun dimère et, chez un XP, sont presque
normaux.
\textbf{13.} $70$ cassures~; au bout d'\qty{1}{h} avec la seule NHEJ,
$70\times
2^{-2} \approx 18$ foyers.
\textbf{14.} $k_{\text{NHEJ}} = \qty{1.39}{h^{-1}}$, $k_{\text{HR}} =
\qty{0.17}{h^{-1}}$~; fraction HR $0.17/1.56 = 0.11$.
\textbf{15.} Demi-vie $\ln 2/1.56 = \qty{27}{min}$~; au bout
d'\qty{1}{h} $70\,e^{-1.56} \approx 15$ cassures.
\textbf{16.} $70\times 0.015\times 0.7 \approx 0.7$ gène détruit par
cellule.
\textbf{17.} $n = 70$~: $2415$ paires $\times 3\times 10^{-4} = 0.72$
translocation. $n = 7$~: $21\times 3\times 10^{-4} = 0.006$. Le nombre
de cassures est $\propto D$, mais le nombre de \emph{paires} de
cassures simultanées est $\propto n^{2} \propto D^{2}$.
\textbf{18.} $D = \qty{7}{h^{-1}}$, $L^{*} = 7/1.39 \approx 5$ cassures
présentes à la fois~: $10$ paires, $0.003$ translocation --- deux cents
fois moins que la même dose en un éclair. Le fractionnement laisse le
tissu sain réparer entre les séances et maintient peu de cassures
simultanées.
\textbf{19.} Réparation intacte~: $e^{-2/1.5} = 0.26$ à \qty{2}{Gy},
$e^{-4} = 0.018$ à \qty{6}{Gy}. Sans NHEJ~: $e^{-4} = 0.018$ et
$e^{-12} = 6\times 10^{-6}$. $D_{0}$ est la dose qui laisse en moyenne
une \omterm{def:b3:dna-repair:lesions}{lésion} létale non réparée par cellule (survie $1/e$).
\textbf{20.} Les \qty{11}{\%} des cassures à deux extrémités de G2 que
la HR aurait réparées fidèlement passent désormais par la NHEJ, et
toute cassure de fourche à une seule extrémité --- que la NHEJ ne sait
pas réparer correctement --- est mal jointe~: au fil des années le
génome accumule délétions, translocations et changements de nombre de
copies, les cicatrices caractéristiques des tumeurs \emph{\omterm{def:b3:dna-repair:dsb}{BRCA}}.
\textbf{21.} $10^{4}\times 0.01 = 100$ cassures à une seule extrémité
par jour. Une cassure à une seule extrémité n'a pas d'extrémité
partenaire~; la NHEJ ne peut que la laisser ou la joindre à une
extrémité étrangère, produisant une translocation ou une perte.
\textbf{22.} Un allèle \emph{BRCA2} fonctionnel produit assez de
protéine pour la recombinaison~: les cellules normales réparent leurs
cassures de fourche et survivent à l'inhibiteur.
\textbf{23.} Un second décalage du cadre, en aval du premier, rétablit
le cadre de lecture et donne une BRCA2 raccourcie mais partiellement
fonctionnelle, si bien que la recombinaison revient et que le
médicament ne tue plus. Se détecte par séquençage de \emph{BRCA2} dans
la tumeur résistante ou dans l'ADN tumoral circulant, et
fonctionnellement par la réapparition des foyers de \omterm{def:b3:dna-repair:dsb}{Rad51} après
irradiation.
\textbf{24.} La mise à mort exige que la \omterm{prop:b3:dna-repair:fidelity}{réparation des mésappariements}
s'attaque aux bases alkylées~; une tumeur de Lynch qui en est dépourvue
y est tolérante --- le médicament ne parvient pas à la tuer (de telles
tumeurs sont traitées autrement, notamment par immunothérapie, à
laquelle leur lourde charge mutationnelle les rend sensibles).
\textbf{25.} Environ \num{92000} dimères à la fourche dans la cellule
XP~; $10^{-10}$ mutation par paire de bases et par division~; la HR
répare \qty{11}{\%} des cassures de G2.
\end{solution}
